\section[Conclusions]{Conclusions}

\begin{frame}{Conclusions}
	Dans ce cours:
	\begin{itemize}
		\item Nous sommes revenus sur l'\textbf{histoire} des réseaux de neurones.
		\pause
		\item Nous avons discuté du premier neurone mathématique, le \textbf{Perceptron}, et de ses limitations.
		\pause
		\item Nous avons vu comment s'affranchir des limitations en combinant \textbf{plusieurs neurones} entre eux grâce à des \textbf{fonctions d'activation non linéaires}.
		\pause
		\item Nous avons alors vu qu'un nouvel \textbf{algorithme} était nécessaire pour l'\textbf{apprentissage}.
	\end{itemize}
\end{frame}

\begin{frame}{Conclusions}
	\begin{itemize}
		\item Nous avons vu comment la \textbf{rétro-propagation du gradient} propage les gradients à travers les différents opérateurs d'un réseau dense (ou complètement connecté).
		\pause
		\item \`A partir des gradients calculés, nous avons vu différents algorithmes pour \textbf{mettre à jour les paramètres} tout au long de l'apprentissage en utilisant des méthodes de batch.
		\pause
		\item Enfin, nous avons discuté des solutions pratiques pour \textbf{initialiser} les paramètres des réseaux de neurones.
	\end{itemize}
\end{frame}